%!TEX root = ../dolgozat.tex
\chapter{Koncepció}
\label{ch:koncepcio}

Ebben a fejezetben azt a szakmai hátteret mutatom be, amelyre a dolgozatban elkészített rendszer
épül. Először a számítógéppel támogatott értékelés általános kérdéseit és a különböző
feladattípusok automatikus pontozását tekintem át, majd részletesebben foglalkozom a programozási
feladatok automatikus értékelésével és a nem megbízható programkód biztonságos futtatásával.
Ezt követően az online vizsgafelügyelet módszereit és az ezekkel kapcsolatos etikai kérdéseket
vizsgálom, majd összehasonlítom a hasonló célú, már létező rendszereket. A fejezet végén
bemutatom a felhasznált technológiákat, és megfogalmazom a rendszerrel szemben támasztott
követelményeket.

\section{Számítógéppel támogatott értékelés}
\label{sec:caa}

A számítógéppel támogatott értékelés (\textit{computer-assisted assessment}, CAA) Bull és
McKenna meghatározása szerint minden olyan értékelési forma, amelyben a számítógép szerepet
játszik a feladatok kiosztásában, a válaszok összegyűjtésében, a pontozásban vagy az
eredmények elemzésében \cite{bull2004blueprint}. A fogalom tehát tágabb, mint az ,,online teszt'':
ide tartozik az is, ha a hallgató papíron dolgozik, de a válaszlapokat gép dolgozza fel, és az is,
ha a hallgató számítógépen ad be egy dolgozatot, amelyet azután az oktató kézzel értékel.
Dolgozatomban a fogalom szűkebb értelmében, a teljes egészében számítógépen zajló vizsgára
használom.

Az értékelés céljuk szerint két nagy csoportba sorolható. A \textit{formatív} értékelés a
tanulási folyamat közben ad visszajelzést a hallgatónak és az oktatónak arról, hol tart a
hallgató, és min kell még dolgoznia; ilyenek a gyakorló feladatok és az önellenőrző tesztek. A
\textit{szummatív} értékelés ezzel szemben egy tanulási szakasz végén méri a tudást, és
jellemzően érdemjegyben vagy pontszámban ölt testet; ilyen a félévközi zárthelyi dolgozat és a
vizsga \cite{bull2004blueprint}. A két forma eltérő követelményeket támaszt egy informatikai
rendszerrel szemben. Formatív értékelésnél a gyors és részletes visszajelzés a legfontosabb, a
csalás elleni védelem kevésbé lényeges, hiszen a hallgatónak magának sem érdeke, hogy
,,becsapja'' a gyakorló feladatot. Szummatív értékelésnél viszont a megbízhatóság, a
megismételhetőség és a vizsga tisztaságának biztosítása kerül előtérbe. Az általam készített
rendszer elsősorban szummatív értékelésre, vagyis vizsgáztatásra készült, de a kódfeladatok
próbafuttatási lehetősége és a nyilvános példatesztek formatív elemeket is visznek bele.

A számítógéppel támogatott értékelés előnyei közül a szakirodalom leggyakrabban a gyorsaságot, a
következetességet és a nagy létszámú csoportok kezelhetőségét emeli ki \cite{bull2004blueprint}.
Egy automatikusan pontozott feladatsor eredménye a beadás pillanatában rendelkezésre áll, és a
pontozás minden hallgatónál ugyanazon szabályok szerint történik, függetlenül attól, hogy az
oktató az első vagy a századik dolgozatot javítja. Az eredmények gépi formában állnak elő, így
könnyen elemezhetők: kiderül, mely feladatok bizonyultak túl könnyűnek vagy túl nehéznek, és
melyik témakörben teljesített gyengén az évfolyam. A hátrányok között leggyakrabban azt
említik, hogy az automatikusan értékelhető feladatok -- különösen a feleletválasztós kérdések --
elsősorban a tényszerű tudást mérik, a magasabb rendű gondolkodási készségek (elemzés,
értékelés, alkotás) mérésére kevésbé alkalmasak, a jó minőségű feladatok összeállítása pedig
jelentős előkészítő munkát igényel \cite{bull2004blueprint}. Saját véleményem szerint ez a
probléma az informatikai képzésben részben feloldható: egy programozási feladat megoldása maga
is alkotó tevékenység, amely ráadásul jól automatizálható módon ellenőrizhető.

\section{Feladattípusok és automatikus pontozásuk}
\label{sec:feladattipusok}

A vizsgafeladatok a válasz kötöttsége szerint zárt, félig nyílt és nyílt feladatokra oszthatók. A
zárt feladatoknál a hallgató előre megadott lehetőségek közül választ; ide tartoznak az
egyválasztós, a több választós és az igaz--hamis kérdések. Ezek pontozása egyértelmű és teljes
mértékben automatizálható, hiszen csak azt kell megvizsgálni, hogy a megjelölt válaszok
egyeznek-e a megoldókulccsal. A félig nyílt feladatoknál a hallgató maga írja be a választ, de
annak formája kötött: egy szám, egy kifejezés vagy egy rövid szöveg. Ezek automatikus értékelése
már nem triviális, hiszen ugyanaz a helyes válasz többféleképpen is leírható (kis- vagy
nagybetűvel, szóközökkel, tizedesvesszővel vagy -ponttal), ezért az összehasonlítás előtt a
válaszokat egységes formára kell hozni, és célszerű több elfogadott választ is megengedni. A
nyílt feladatok, például a kifejtős kérdések, a mai eszközökkel megbízhatóan nem értékelhetők
automatikusan, ezeknél a számítógép szerepe a válaszok összegyűjtésére és a kézi javítás
támogatására korlátozódik.

Külön figyelmet érdemel a több helyes választ tartalmazó feladatok pontozása. Ha a rendszer csak a
tökéletes választ fogadja el, akkor egyetlen elnézett jelölés is nulla pontot eredményez, ami a
hallgató tudását nem tükrözi hűen. Ha viszont minden helyes jelölésért részpont jár, és a hibás
jelölésekért nem jár levonás, akkor a hallgató egyszerűen az összes lehetőséget megjelölheti, és
így biztosan megkapja a teljes pontszámot. A gyakorlatban ezért elterjedt megoldás, hogy a helyes
jelölések részpontot érnek, a hibás jelölések pedig ugyanilyen mértékben csökkentik a pontszámot,
de a feladatra adott pontszám nem lehet negatív. A rendszeremben ezt a pontozási módot
valósítottam meg (lásd \aref{sec:pontozas}. szakaszt), mert véleményem szerint ez egyensúlyt teremt
a részleges tudás elismerése és a találgatás visszaszorítása között.

\section{Programozási feladatok automatikus értékelése}
\label{sec:apas}

A programozási feladatok automatikus értékelése egyidős a programozásoktatással. Már 1960-ban
megjelent egy cikk, amely egy olyan programról számolt be, amely a hallgatók által lyukkártyán
beadott programokat futtatta le és értékelte \cite{hollingsworth1960automatic}. Douce és szerzőtársai
áttekintő tanulmányukban a terület fejlődését három generációra osztották
\cite{douce2005automatic}. Az első generáció rendszerei egy-egy konkrét környezethez kötött,
egyedi megoldások voltak. A második generáció általánosabb, parancssori és grafikus
eszközöket hozott, amelyek már több feladattípust és programozási nyelvet támogattak. A harmadik
generációt a webes rendszerek jelentik, amelyekben a hallgató böngészőből adja be a megoldását, és
az eredményt is ott kapja meg.

Ihantola és szerzőtársai a 2006 és 2010 között megjelent rendszereket vizsgálva arra a
következtetésre jutottak, hogy bár sok rendszer készült, ezek jelentős része egy-egy intézmény
belső eszköze maradt, és a rendszerek közötti legfontosabb különbségek a támogatott programozási
nyelvekben, a tesztek megadásának módjában, az újrabeküldés szabályaiban, a kézi értékelés
lehetőségében és a futtatás biztonságában mutatkoznak meg \cite{ihantola2010review}. Különösen az
utolsó szempontot emelték ki: egy automatikus értékelő rendszernek úgy kell futtatnia a
hallgatók programjait, hogy azok ne veszélyeztethessék sem a szervert, sem más hallgatók
megoldásait. Paiva és szerzőtársai egy jóval újabb, 2022-es áttekintésükben megállapították, hogy
az automatikus értékelés mára a számítástechnika-oktatás széles körben használt eszközévé vált,
és a kutatások súlypontja egyre inkább a hallgatónak adott visszajelzés minőségére helyeződik át
\cite{paiva2022automated}.

Az értékelés módszerei alapvetően két csoportba sorolhatók \cite{paiva2022automated}. A
\textit{dinamikus} elemzés a programot ténylegesen lefuttatja, és a viselkedését vizsgálja;
ennek legelterjedtebb formája a tesztelésen alapuló értékelés, amelyben a program előre megadott
bemenetekre adott kimenetét hasonlítjuk össze az elvárt kimenettel. A \textit{statikus} elemzés
a programot futtatás nélkül, a forráskód alapján vizsgálja, például a kódolási stílus, a
bonyolultság vagy bizonyos tiltott nyelvi elemek használata szempontjából. Rendszeremben a
dinamikus, tesztelésen alapuló értékelést valósítottam meg, mert ez közvetlenül azt méri, amit a
vizsgafeladat számon kér: hogy a program helyesen oldja-e meg a feladatot.

A tesztelésen alapuló értékelés egyik fontos tervezési kérdése, hogy a hallgató mennyit lásson a
tesztesetekből. Ha minden tesztesetet ismer, akkor fennáll a veszély, hogy a programját a konkrét
tesztekre ,,szabja'', például az elvárt kimeneteket egyszerűen kiírja, ahelyett hogy általános
megoldást készítene. Ha viszont egyet sem ismer, akkor nem tudja ellenőrizni, hogy jól értette-e
a feladatot, és például a kimenet formátumán (egy fölösleges szóközön vagy sortörésen) is
elbukhat. A gyakorlatban elterjedt kompromisszum a tesztesetek két csoportra bontása: a
nyilvános példatesztek a feladat megértését segítik, a rejtett tesztek pedig a megoldás
általánosságát ellenőrzik. Ezt a megoldást alkalmaztam a rendszeremben is.

\section{Nem megbízható programkód biztonságos futtatása}
\label{sec:sandbox_elmelet}

Az automatikus értékelő rendszer a hallgatók programjait futtatja, vagyis olyan kódot, amelynek
tartalmáról a futtatás előtt semmit sem tudunk. Egy ilyen programot -- a hibás és a szándékosan
ártó programokat egyaránt -- \textit{nem megbízható kódnak} nevezünk. A lehetséges fenyegetések
közé tartozik az erőforrások kimerítése (végtelen ciklus, korlátlan memóriafoglalás, rengeteg
folyamat indítása, óriási mennyiségű kimenet előállítása), a szerver fájljainak olvasása vagy
módosítása (például más hallgatók megoldásainak vagy a tesztesetek elvárt kimeneteinek
megszerzése), hálózati kapcsolatok kezdeményezése, valamint az értékelés eredményének
meghamisítása. A futtatókörnyezetet ezért úgy kell kialakítani, hogy a program csak a
számára kijelölt, szigorúan korlátozott erőforrásokhoz férhessen hozzá. Az ilyen elszigetelt
futtatókörnyezetet nevezzük homokozónak (\textit{sandbox}).

Az elszigetelésre több megoldás is létezik. A legegyszerűbb az operációs rendszer
folyamatszintű korlátozásainak használata (időkorlát, memóriakorlát, jogosultságok nélküli
felhasználó), ez azonban önmagában nem akadályozza meg, hogy a program lássa a szerver
fájlrendszerét és hálózatát. A másik véglet a teljes virtuális gép, amely erős elszigetelést nyújt,
de az indítása lassú és sok erőforrást igényel, ami egy vizsga során egyszerre beküldött több
száz megoldás esetén komoly korlátot jelent. A kettő között helyezkednek el a konténerek.

A Docker konténerek a Linux kernel két szolgáltatására épülnek \cite{merkel2014docker}. A
\textit{névterek} (\textit{namespaces}) segítségével a konténerben futó folyamat saját, elkülönített
nézetet kap a rendszerről: saját folyamatlistát, fájlrendszert, hálózati csatolókat és
felhasználókat lát, így a gazdagép többi folyamatához és fájljához nem fér hozzá. A
\textit{vezérlőcsoportok} (\textit{control groups}, cgroups) pedig lehetővé teszik a konténer által
felhasználható erőforrások -- memória, processzoridő, folyamatok száma -- korlátozását
\cite{docker_docs}. A konténer nem tartalmaz saját operációsrendszer-kernelt, hanem a gazdagépét
használja, ezért a virtuális gépeknél lényegesen gyorsabban, jellemzően egy másodpercen belül
elindul \cite{merkel2014docker}. Ennek az ára, hogy az elszigetelés gyengébb: egy kernelszintű
biztonsági rés esetén a konténerből elvileg ki lehet törni. Ezért a konténer jogosultságait is a
lehető legszűkebbre kell szabni: le kell tiltani a hálózatot, csak olvashatóvá kell tenni a
fájlrendszert, el kell dobni a Linux különleges jogosultságait (\textit{capabilities}), és a
programot nem rendszergazdaként kell futtatni. Az általam alkalmazott beállításokat \aref{tab:sandbox}.
táblázat foglalja össze.

\section{Online vizsgafelügyelet és akadémiai integritás}
\label{sec:proctoring}

Az online vizsgák egyik legtöbbet vitatott kérdése a vizsgázók felügyelete
(\textit{proctoring}). A felügyeleti megoldások három fő típusa különböztethető meg
\cite{nigam2021systematic}. Az \textit{élő felügyelet} során egy emberi felügyelő webkamerán
keresztül valós időben figyeli a vizsgázókat. A \textit{rögzített felügyelet} esetén a vizsga
alatt kép- és hangfelvétel készül, amelyet később egy felügyelő átnéz. Az \textit{automatizált
felügyelet} pedig mesterséges intelligenciára épülő módszerekkel elemzi a vizsgázó viselkedését,
és csak a gyanús eseteket jelzi emberi ellenőrzésre. Nigam és szerzőtársai rendszerező
áttekintésükben azt találták, hogy a COVID-19 járvány hatására az automatizált felügyeleti
rendszerek iránti kereslet jelentősen megnőtt, ugyanakkor ezek a rendszerek komoly kihívásokkal
küzdenek, többek között az adatvédelem, a téves riasztások és a technikai feltételek terén
\cite{nigam2021systematic}.

Az automatizált felügyelet egyik sokat hivatkozott példája Atoum és szerzőtársai rendszere,
amely két kamera (egy webkamera és egy viselhető kamera) és egy mikrofon jeleit dolgozza fel
\cite{atoum2017automated}. A rendszer több alapvető viselkedési jellemzőt figyel, például a
vizsgázó személyazonosságát, a tekintet irányát, a beszédet, a mobiltelefon jelenlétét és a
számítógépen aktív ablakot, majd ezek alapján osztályozza, hogy történt-e csalás. Érdemes
kiemelni, hogy az aktív ablak figyelése, vagyis annak észlelése, hogy a vizsgázó másik
alkalmazásra vált, ebben a komplex rendszerben is külön jellemzőként szerepel.

A felügyeleti rendszerek használata azonban komoly etikai kérdéseket vet fel. Coghlan, Miller és
Paterson részletesen elemezték az online egyetemi vizsgák felügyeletének etikai vonatkozásait
\cite{coghlan2021proctor}. Megállapításaik szerint a webkamerás és mesterséges intelligenciára épülő
felügyelet a hallgatók magánszféráját érinti (hiszen a felügyelet a hallgató otthonába ,,lát
be''), a téves riasztások igazságtalan gyanúba keverhetik a vizsgázókat, a rendszerek működése
pedig sokszor nem átlátható. Ugyanakkor nem vetik el teljesen a felügyelet alkalmazását: úgy
érvelnek, hogy bizonyos esetekben indokolt lehet, ha az intézmény átláthatóan tájékoztatja a
hallgatókat, a beavatkozás mértéke arányos a céllal, és az automatikus jelzések alapján
emberi döntés születik \cite{coghlan2021proctor}.

Ezek alapján a saját rendszeremben egy szándékosan szerény, a magánszférát nem sértő megoldást
választottam: a rendszer kizárólag azt figyeli, hogy a vizsgázó böngészőablaka elveszíti-e a
fókuszt, illetve elrejti-e a vizsgaoldalt, és ezt valós időben jelzi az oktatónak. Kamera- és
hangfelvétel nem készül, a rendszer nem hoz automatikus döntést, csupán információt ad az
oktatónak, aki a jelzéseket a vizsga körülményeinek ismeretében értékelheti. Tisztában vagyok vele,
hogy ez a módszer nem nyújt teljes védelmet, hiszen egy második eszköz használatát nem képes
észlelni; véleményem szerint azonban egy egyetemi vizsgán a visszatartó hatás és az oktató
tájékoztatása a magánszféra tiszteletben tartásával együtt jobb egyensúlyt jelent, mint a
hallgatók folyamatos megfigyelése.

\section{Hasonló alkalmazások}
\label{sec:hasonlo}

A rendszer megtervezése előtt megvizsgáltam néhány elterjedt, hasonló célú alkalmazást. Ezek közül
hármat mutatok be részletesebben, mert jól reprezentálják a lehetséges megközelítéseket.

A \textbf{Moodle} nyílt forráskódú tanulásirányítási rendszer (LMS) a felsőoktatásban az egyik
legelterjedtebb e-learning platform, amelyet hazánkban is számos egyetem használ. Teszt (\textit{Quiz})
modulja sokféle kérdéstípust támogat, többek között feleletválasztós, igaz--hamis, rövid
válaszos, numerikus és esszé típusú kérdéseket, valamint részletes beállításokat kínál az
időkorlátra, a próbálkozások számára és a visszajelzésre \cite{moodle_quiz_docs}. A tesztmodul
önmagában nem futtat programkódot, és valós idejű felügyeletet sem nyújt; a vizsgakörnyezet
korlátozására a Safe Exam Browser nevű külön telepítendő, zárolt böngészővel integrálható
\cite{seb_site}.

A \textbf{CodeRunner} a Moodle egy bővítménye, amely új kérdéstípusként programozási feladatokat
ad a tesztmodulhoz \cite{coderunner_site}. A hallgató által beküldött kódot egy különálló,
homokozóként működő szerveren futtatja le tesztesetekkel, és támogatja a hallgató elől elrejtett
teszteseteket is. Mivel a Moodle tesztmodulján belül működik, a többi kérdéstípussal együtt
használható, így vegyes vizsgák is összeállíthatók vele. Hátránya, hogy használatához egy teljes
Moodle-rendszer és egy külön kódfuttató szerver üzemeltetése szükséges, ami egy kisebb oktatói
csoport számára jelentős üzemeltetési terhet jelent.

A \textbf{DOMjudge} programozási versenyek (például az ICPC) lebonyolítására készült nyílt
forráskódú rendszer \cite{domjudge_site}. A beküldött programokat elszigetelt környezetben futtatja,
és rejtett tesztadatok alapján dönti el, hogy a megoldás helyes-e. A versenyekhez igazodva
azonban csak programozási feladatokat kezel, az értékelés jellemzően ,,elfogadva vagy elutasítva''
jellegű, részpontozásra és kifejtős feladatokra nincs lehetőség, így egy egyetemi vizsga teljes
lebonyolítására kevésbé alkalmas.

\begin{table}[H]
  \centering
  \caption{A vizsgált rendszerek összehasonlítása}
  \label{tab:osszehasonlitas}
  \small
  \begin{tabularx}{\textwidth}{@{}>{\raggedright\arraybackslash}X>{\centering\arraybackslash}p{2.5cm}>{\centering\arraybackslash}p{2.6cm}>{\centering\arraybackslash}p{2.0cm}>{\centering\arraybackslash}p{2.0cm}@{}}
    \toprule
    Szempont & Moodle Quiz & Moodle + CodeRunner & DOMjudge & Saját rendszer \\
    \midrule
    Vegyes feladattípusok        & igen & igen & nem & igen \\
    Programkód futtatása         & nem & igen & igen & igen \\
    Rejtett tesztesetek          & -- & igen & igen & igen \\
    Részpontozás kódfeladatnál   & -- & igen & nem & igen \\
    Kézi javítás, visszajelzés   & igen & igen & részben & igen \\
    Ablakelhagyás valós idejű jelzése & nem (zárolt böngészővel korlátozható) & nem & nem & igen \\
    Üzemeltetési igény           & nagy & nagy & közepes & kicsi \\
    \bottomrule
  \end{tabularx}
\end{table}

A rendszerek főbb jellemzőit \aref{tab:osszehasonlitas}. táblázat foglalja össze. Az összevetés
alapján úgy látom, hogy a kódfuttatást és a vegyes feladattípusokat a Moodle és a CodeRunner
együttese kínálja a legteljesebb formában, ennek azonban jelentős üzemeltetési ára van, a valós
idejű felügyelet pedig ott sem megoldott. Az általam készített rendszer nem kíván versenyezni egy
teljes tanulásirányítási rendszerrel, hanem azt mutatja meg, hogy a vizsgáztatás e három
összetevője -- a vegyes feladattípusok, a biztonságos kódfuttatás és a valós idejű felügyelet --
egy könnyen telepíthető, átlátható felépítésű alkalmazásban is egyesíthető.

\section{Felhasznált technológiák}
\label{sec:technologiak}

A technológiák kiválasztásánál fontos szempont volt, hogy elterjedt, jól dokumentált, nyílt
forráskódú eszközöket használjak, amelyek a webes alkalmazásfejlesztés jelenlegi gyakorlatát
képviselik, és amelyekkel egy kisebb csapat is hatékonyan tud dolgozni.

\subsection{REST architektúra}

A kliens és a szerver közötti kommunikáció a REST (\textit{Representational State Transfer})
architekturális stílust követi, amelyet Roy Fielding fogalmazott meg doktori értekezésében
\cite{fielding2000architectural}. A REST egy sor megszorítást határoz meg: a kliens és a szerver
elkülönülését, az állapotmentes kommunikációt (minden kérés tartalmazza a feldolgozásához
szükséges összes információt), a gyorsítótárazhatóságot, az egységes interfészt, valamint a
rétegzett felépítést \cite{fielding2000architectural}. A gyakorlatban ez azt jelenti, hogy a
rendszer erőforrásai (vizsgák, feladatok, kitöltések, válaszok) URL-ekkel azonosíthatók, a rajtuk
végzett műveleteket pedig a HTTP metódusai (GET, POST, PUT, PATCH, DELETE) fejezik ki. Az
állapotmentesség különösen előnyös a hitelesítés szempontjából, mert lehetővé teszi a
token alapú azonosítást, ahol a szervernek nem kell munkamenetet tárolnia.

\subsection{Szerveroldal: Python, FastAPI, SQLAlchemy, PostgreSQL}

A szerveroldalt Python nyelven, a FastAPI keretrendszerrel készítettem el \cite{fastapi_docs}. A
FastAPI az ASGI (\textit{Asynchronous Server Gateway Interface}) szabványra épül, így szinkron és
aszinkron kéréskezelőket egyaránt támogat, ami a valós idejű kommunikációhoz használt Socket.IO
szerver beépítését is megkönnyítette. A keretrendszer a Python típusannotációit használja fel a
kérések paramétereinek feldolgozására, a Pydantic könyvtár segítségével pedig a be- és kimenő
adatok automatikus ellenőrzésére és átalakítására \cite{pydantic_docs}. Ennek nagy előnye, hogy a
validációs szabályokat egyszer, deklaratív módon kell megfogalmazni, és a keretrendszer ezekből
automatikusan interaktív API-dokumentációt (OpenAPI/Swagger) is előállít, ami a fejlesztés és a
tesztelés során sokat segített. Alternatívaként a Django és a Flask keretrendszer jöhetett
volna szóba; a Django egy teljes, ,,mindent tartalmazó'' keretrendszer, amely egy REST API-hoz
fölöslegesen nehézkes, a Flaskhoz pedig a validációt és a dokumentációt külön könyvtárakkal
kellett volna megoldani.

Az adatokat PostgreSQL relációs adatbázisban tárolom \cite{postgresql_docs}, amely nyílt forráskódú,
megbízható, és a tranzakciókezelés, az egyedi megszorítások és a JSON típusú oszlopok
támogatása miatt jól illeszkedik a feladathoz. Az adatbázist az SQLAlchemy objektum-relációs
leképező (ORM) könyvtáron keresztül érem el \cite{sqlalchemy_docs}, amely a táblákat Python
osztályokként, a rekordokat objektumokként kezeli. Ennek köszönhetően a kód nagy része
adatbázis-függetlenné vált: fejlesztés és tesztelés közben ugyanaz a kód SQLite adatbázissal is
működik.

\subsection{Hitelesítés: JWT és bcrypt}

A felhasználók azonosítására JSON Web Tokent (JWT) használok. A JWT egy tömör, URL-biztos
formátum, amely egy fejlécből, egy adatrészből (\textit{payload}) és egy digitális aláírásból áll,
ezeket pont választja el egymástól \cite{rfc7519}. Az adatrész úgynevezett állításokat
(\textit{claims}) tartalmaz, például a felhasználó azonosítóját (\texttt{sub}) és a token lejárati
idejét (\texttt{exp}). Mivel a tokent a szerver írja alá, annak tartalma a kliens által nem
módosítható észrevétlenül, így a szerver minden kérésnél adatbázis-alapú munkamenet nélkül
ellenőrizheti a felhasználó személyazonosságát.

A jelszavak tárolására a bcrypt algoritmust választottam, amelyet Provos és Mazières kifejezetten
jelszavak hash-elésére terveztek \cite{provos1999bcrypt}. A bcrypt minden jelszóhoz véletlen sót
(\textit{salt}) használ, így két azonos jelszó hash-e is különböző, és egy állítható
költségparaméterrel szabályozható, hogy egy hash kiszámítása mennyi időt vegyen igénybe. A
szerzők célja az volt, hogy az algoritmus ,,együtt tudjon nőni'' a hardverek teljesítményével:
a költségparaméter emelésével a nyers erővel történő jelszótörés a jövőben is költséges
maradhat \cite{provos1999bcrypt}.

\subsection{Valós idejű kommunikáció: WebSocket és Socket.IO}

A HTTP protokoll kérés--válasz alapú: a kommunikációt mindig a kliens kezdeményezi. A valós idejű
felügyelethez azonban arra van szükség, hogy a szerver saját kezdeményezésből is üzenhessen az
oktató kliensének, amikor egy hallgató elhagyja a vizsgaablakot. Erre szolgál a WebSocket
protokoll, amely egy HTTP-kérésként induló kézfogás (\textit{handshake}) után egy állandó,
kétirányú (\textit{full-duplex}) kapcsolatot épít ki a kliens és a szerver között \cite{rfc6455}.

A WebSocketet nem közvetlenül, hanem a Socket.IO könyvtáron keresztül használom
\cite{socketio_docs}. A Socket.IO a WebSocket fölé egy eseményalapú programozási felületet
helyez, automatikusan újracsatlakozik a kapcsolat megszakadása esetén, és ha a WebSocket nem
érhető el (például egy proxy kiszűri), HTTP-alapú hosszú lekérdezésre (\textit{long polling})
vált. Különösen hasznos a szobák (\textit{rooms}) fogalma, amellyel az üzenetek egy kijelölt
klienscsoportnak, esetemben egy adott vizsga résztvevőinek küldhetők el.

\subsection{Kliensoldal: React, Next.js, Tailwind CSS}

A felhasználói felületet a React könyvtárra épülő Next.js keretrendszerrel készítettem
\cite{react_docs,nextjs_docs}. A React komponensalapú megközelítése lehetővé teszi, hogy a felület
újrahasznosítható építőelemekből álljon össze; a vizsgaoldalon például ugyanaz a válaszbeviteli
komponens jeleníti meg mind a hat feladattípust. A React \textit{hook} mechanizmusa lehetővé teszi,
hogy egy összetett, állapottal rendelkező logikát -- mint a felügyeleti kapcsolat kezelését --
egy önálló, újrahasznosítható függvénybe szervezzek ki. A Next.js ehhez fájlrendszer-alapú
útvonalkezelést (az App Routert) és egyszerű fejlesztői környezetet ad. A TypeScript nyelv
használata a szerver által küldött adatok típusainak leírásával sok hibát már fordítási időben
kiszűrt. A megjelenéshez a Tailwind CSS segédosztály-alapú keretrendszert használtam
\cite{tailwind_docs}, amellyel a stílusok közvetlenül a komponensekben adhatók meg, így nem kellett
külön stíluslapokat karbantartani.

\subsection{Konténerizáció: Docker}

A Docker \cite{merkel2014docker,docker_docs} a rendszerben kettős szerepet tölt be. Egyrészt a
fejlesztési környezet része: a PostgreSQL adatbázis egy Docker Compose leírással egyetlen
paranccsal elindítható, így a rendszer telepítése nem igényli az adatbázis-kezelő külön
telepítését. Másrészt -- és ez a fontosabb szerep -- ez biztosítja a hallgatók programjainak
elszigetelt futtatását, amelynek elméleti hátterét \aref{sec:sandbox_elmelet}. szakaszban
ismertettem. A szerveroldali alkalmazás a Docker Python nyelvű programozói felületén
(SDK-ján) keresztül indítja és állítja le a konténereket.

\section{Követelményspecifikáció}
\label{sec:kovetelmenyek}

A követelmények megfogalmazásánál a szoftvertechnológia szokásos felosztását követtem, amely
megkülönbözteti a rendszer által nyújtandó szolgáltatásokat leíró funkcionális, valamint a
rendszer működésének minőségi jellemzőit és korlátait leíró nem funkcionális követelményeket
\cite{sommerville2016software}.

\subsection{Szerepkörök és funkcionális követelmények}

A rendszernek két szerepköre van. A \textbf{hallgató} vizsgákat tölt ki, az \textbf{oktató} pedig
vizsgákat hoz létre, felügyeli és értékeli azokat. A regisztráció alapértelmezés szerint hallgatói
fiókot hoz létre; oktatói fiók csak egy meghívókód ismeretében regisztrálható. A szerepkörökhöz
tartozó főbb funkcionális követelményeket \aref{tab:funkcionalis}. táblázat foglalja össze.

\begin{longtable}{@{}lp{1.9cm}p{10.9cm}@{}}
  \caption{Funkcionális követelmények}\label{tab:funkcionalis}\\
  \toprule
  Azon. & Szerepkör & Követelmény \\
  \midrule
  \endfirsthead
  \toprule
  Azon. & Szerepkör & Követelmény \\
  \midrule
  \endhead
  F1  & mindkettő & Regisztráció e-mail címmel és jelszóval, bejelentkezés, kijelentkezés. \\
  F2  & oktató    & Vizsga létrehozása címmel, leírással, időablakkal (kezdés és vége) és kitöltési időtartammal; vizsga törlése. \\
  F3  & oktató    & Feladatok hozzáadása és törlése hat típusban: egyválasztós, több választós, igaz--hamis, rövid szöveges, kifejtős és programozási feladat. \\
  F4  & oktató    & Programozási feladathoz tesztesetek megadása (bemenet, elvárt kimenet), mindegyikről eldöntve, hogy nyilvános példateszt vagy rejtett. \\
  F5  & hallgató  & Az elérhető vizsgák listájának megtekintése, vizsga elindítása az időablakon belül, vizsgánként egyszer. \\
  F6  & hallgató  & Feladatok kitöltése visszaszámlálóval, a válaszok automatikus mentése. \\
  F7  & hallgató  & Programkód próbafuttatása saját bemenettel vagy a nyilvános példatesztekkel; a kód beküldése értékelésre. \\
  F8  & hallgató  & Vizsga beadása; az idő lejártakor automatikus beadás. \\
  F9  & hallgató  & Az eredmény megtekintése a beadás után, feladatonkénti pontszámmal és oktatói visszajelzéssel. \\
  F10 & rendszer  & A zárt, a rövid szöveges és a programozási feladatok automatikus pontozása. \\
  F11 & oktató    & A vizsga élő felügyelete: a hallgatók jelenlétének és ablakelhagyásainak valós idejű megjelenítése. \\
  F12 & oktató    & Eredmények megtekintése, kifejtős feladatok pontozása, automatikus pontszám felülírása, szöveges visszajelzés. \\
  F13 & oktató    & Eredmények exportálása táblázatkezelővel megnyitható (CSV) formátumban; sikertelen kódértékelések újrafuttatása. \\
  \bottomrule
\end{longtable}

A rendszer működését egy tipikus forgatókönyv szemlélteti. Az oktató létrehoz egy vizsgát, amely
egy adott napon 10 és 12 óra között indítható el, és 60 perc alatt kell kitölteni. A vizsgához
hozzáad egy feleletválasztós, egy rövid válaszos, egy kifejtős és egy programozási feladatot, az
utóbbihoz két nyilvános és három rejtett tesztesettel. A vizsga napján a hallgató 10:20-kor
elindítja a vizsgát, így a határideje 11:20 lesz. Kitölti a feladatokat, a programját a
példatesztekkel kipróbálja, majd beküldi. Közben egyszer átvált egy másik böngészőfülre, amit
az oktató az élő monitoron azonnal lát. A beadás után a zárt feladatok és a programozási
feladat pontszáma azonnal elkészül, a kifejtős választ pedig az oktató értékeli, és szöveges
visszajelzést fűz hozzá, amelyet a hallgató a saját felületén lát.

\subsection{Nem funkcionális követelmények}

A nem funkcionális követelmények közül a biztonsági követelményeket tartottam a
legfontosabbnak. A webalkalmazások leggyakoribb biztonsági kockázatait összegző OWASP Top Ten
lista első helyén a hozzáférés-szabályozás hibái szerepelnek \cite{owasp_top10}, ezért külön
hangsúlyt fektettem arra, hogy minden végpont ellenőrizze a felhasználó szerepkörét és az
adott erőforráshoz való jogosultságát. A rendszerrel szemben a következő nem funkcionális
követelményeket fogalmaztam meg:
\begin{itemize}
  \item \textbf{Bizalmasság:} a megoldókulcsok, a rejtett tesztesetek és a többi hallgató
  válaszai semmilyen módon nem juthatnak el a hallgatóhoz; a pontszám csak a beadás után
  látható.
  \item \textbf{Elszigetelés:} a hallgatók programjai nem férhetnek hozzá a szerver
  fájljaihoz és a hálózathoz, és nem meríthetik ki a szerver erőforrásait.
  \item \textbf{Megbízhatóság:} a hallgató munkája ne vesszen el hálózati hiba, oldalfrissítés
  vagy a böngésző bezárása esetén sem; rendszerhiba miatt a hallgató ne kapjon nulla pontot.
  \item \textbf{Igazságosság:} a határidőt a szerver érvényesítse, ne a hallgató számítógépének
  órája; mindenki ugyanannyi időt kapjon.
  \item \textbf{Válaszidő:} egy programozási feladat értékelése néhány másodpercen belül
  elkészüljön, a felügyeleti jelzés pedig érzékelhető késés nélkül jelenjen meg.
  \item \textbf{Hordozhatóság és telepíthetőség:} a rendszer Windows, Linux és macOS
  operációs rendszeren is futtatható legyen, minimális telepítési lépéssel.
\end{itemize}
